Let $\Delta V_0$ be the voltage of one battery. The reference bulb gets $\Delta V_0$; the two batteries in series give $2\,\Delta V_0$.

A is in series with the pair B\,$\|$\,C: the whole current flows through A and then splits between B and C. At the same voltage the pair lets more current through than one bulb, because it offers two paths. In a series connection the part that lets less current through takes the larger share of the voltage (voltage divider). So
\begin{align}
 \Delta V_A > \Delta V_0 \qquad\text{and}\qquad \Delta V_B = \Delta V_C = 2\,\Delta V_0 - \Delta V_A < \Delta V_0
\end{align}

\begin{itemize}
\item A has more voltage than $L_0$, and so also more current (the current of a bulb grows with its voltage): $P_A = \Delta V_A\cdot I_A > P_0$.
\item B and C have less voltage than $L_0$ and so less current: $P_B = P_C < P_0$. They are identical and in parallel, so they are equally bright, each with half the current of A.
\end{itemize}

The ranking is
\begin{align}
 \result{P_A > P_0 > P_B = P_C}
\end{align}
A is brighter than the reference bulb, B and C are less bright, and equally bright.

A check with bulbs of constant resistance $R$ (real bulbs are not ohmic, but the ranking is the same): the pair has $R/2$, so A gets $\frac{2}{3}$ of $2\,\Delta V_0$, that is $\frac{4}{3}\,\Delta V_0$, and the pair $\frac{2}{3}\,\Delta V_0$. With $P = \Delta V^2/R$ this gives $P_A = \frac{16}{9}\,P_0$ and $P_B = P_C = \frac{4}{9}\,P_0$.
